\documentclass[10pt,a4paper]{article}
\usepackage{amsmath,amssymb,amsthm,graphicx}
\usepackage{hyperref}
\usepackage{latexsym}
\usepackage{amssymb}
\usepackage{amsthm}
\usepackage{ulem}
\usepackage{enumitem}
\setlist{itemsep=0.4pt,topsep=0pt}
\parindent 0pt
\parindent 0pt
\parskip 12pt
\newcommand{\HCF}{\mathrm{hcf}}
\newcommand{\ZZ}{\mathbb{Z}}
\newcommand{\rk}{\mathrm{rk}}
\DeclareMathOperator{\len}{length}
\DeclareMathOperator{\bl}{b}
\DeclareMathOperator{\cl}{cl}
\DeclareMathOperator{\di}{d}
\DeclareMathOperator{\Cl}{Cl}
\DeclareMathOperator{\acl}{acl}
\DeclareMathOperator{\tp}{tp}
\DeclareMathOperator{\Th}{Th}
\newcommand{\Cor}{\mathcal{C}_\mathrm{or}}
\newcommand{\Cun}{\mathcal{C}_\mathrm{un}}
\newcommand{\mons}{\mathbb{M}}
\newcommand{\strong}{\mathrm{strong}}
\newcommand{\axioms}{\Delta_{\mathrm{FSTS}}}
\newtheoremstyle{note}% hnamei
{8pt}% hSpace abovei
{14pt}% hSpace belowi
{\normalfont}% hBody fonti
{}% hIndent amounti1
{\bfseries}% hTheorem head fonti
{ }% hPunctuation after theorem headi
{.5em}% hSpace after theorem headi2
{}% hTheorem head spec (can be left empty, meaning ‘normal’)i
\theoremstyle{note}
\newtheorem{thm}{Theorem}[section]
%\newtheorem{theorem}{Theorem}
\newtheorem{lemma}[thm]{Lemma} 
\newtheorem{df}[thm]{Definition}
\newtheorem{rmk}[thm]{Remark} 
\newtheorem{prop}[thm]{Proposition}
\newtheorem{fact}[thm]{Fact}
\newtheorem{cor}[thm]{Corollary}
\newtheorem{conj}[thm]{Conjecture}
\newtheorem{ex}[thm]{Example}
\newtheorem{qu}[thm]{Question}
\newtheorem*{notation}{Notation} 
\title{An isomorphism of partial subsystems over elementary subsystems is partially elementary}
\author{Silvia Barbina, Enrique Casanovas}
\date{arxiv:2411.13723v3; Creative Commons Attribution 4.0 International}
\begin{document}
\maketitle
\noindent\textit{Editorial note.} The counted statement is the author's environment at source0.tex lines 1001-1002 and its proof environment opens at line 1003; the pairing is by adjacency and the retained proof range is the whole environment.\par\medskip
\noindent\textit{Editorial note.} The preamble is the author's own package lines, macro block and theorem-environment declarations (source lines 4-65); the author's optioned amsmath line and the fleqn-dependent layout line are not retained.\par\medskip
\noindent\textit{Editorial note.} The source carries no label-number dump and no counter values are imposed, so the printed numbering is the wrapper's default; identity is the author label, or the pinned source position where the author gave none.\par\medskip
\noindent\textit{Editorial note.} Metadata provenance: the source record of the pinned copy and the licence evidence file are both present; the licence is CC BY 4.0 with status supported, and the card's source\_version matches the value recorded in that evidence file.\par\medskip
\noindent\textit{Editorial note.} r002 note: the retained May\_26\_2024\_2 context refers to Definition def5\_1, so r002 adds exactly that definition as a further context and changes nothing else.\par\medskip
\noindent\textit{Editorial scope.} One original result of the pinned author manuscript, counted once: isomorphisms of partial subsystems over elementary subsystems are partially elementary (author label elementary).\par\medskip
\noindent\textit{Allowed standard prerequisites (editorial).}\begin{itemize}
\item Partial elementary map. A map between subsets is partial elementary when it preserves the truth of formulas on the tuples where it is defined.. It is the conclusion of the statement.
\item Elementary subsystem. A is elementary in M when the formulas with parameters in A agree.. It is the hypothesis on both sides.
\item Saturated elementary extension. Every structure has a saturated elementary extension of prescribed large cardinality.. It is the reduction used in the countable case.
\end{itemize}\medskip
	\begin{df} \label{def5_1} For $A, B \in \Cun$,
$$A \leq B \  \Leftrightarrow \ A\subseteq B \mbox{ and there is an HF-ordering of $B$ where $A$ is an initial segment}. $$
When $A \leq B$, we say that $A$ is \textbf{strong} in $B$.
\end{df}

\begin{lemma}\label{Feb_23_2024_7} If $<$ is an HF-ordering of $B$  and $A=\cl_<(A)\subseteq B$, then  $<$ is an HF-ordering of $B$ over $A$. 
\end{lemma}

\begin{lemma} \label{May_26_2024_2} Let $M$ be an unconfined STS.  If $A\leq M$ (as in Definition~\ref{def5_1}), then  $\langle A\rangle \leq M$ and there is an HF-ordering of $\langle A\rangle$ over $A$ whose final segment $\langle A\rangle \smallsetminus A$ is well-ordered. If $N$ is an unconfined STS and $B\leq N$ and $f:A\cong B$ is an isomorphism of partial STSs, then $f$ can be extended to an isomorphism $f^\prime :\langle A\rangle \cong \langle B \rangle$.
\end{lemma}

\begin{thm} \label{May_26_2024_4} Let $M,N$ be uncountable saturated unconfined STSs. The set of isomorphisms between countable substructures $A\subseteq M$ and $B\subseteq N$  (in particular,  $A,B$ are STSs)  such that $A\leq M$ and $B\leq N$ is a back-and-forth system between $M$ and $N$.
\end{thm}	

\begin{cor}\label{elementary} Let $M,N$ be infinite unconfined STS. If $A\leq M$, $B\leq N$ and $f:A\rightarrow B$ is an isomorphism of partial STS, then $f$ is a partial elementary map from $M$ to $N$.
	\end{cor}

	\begin{proof} First consider the case where $A$ and $B$ are countable.  By Lemma~\ref{May_26_2024_2}, $\langle A\rangle \leq M$,  $\langle B\rangle \leq N$  and $f$ extends to some isomorphism  $f^\prime:\langle A\rangle \cong \langle B\rangle $.
		 We can find uncountable saturated (or special $\omega_1$-saturated) elementary extensions  $M^\prime\succeq M$ and $N^\prime\succeq N$. Notice that  $A\leq M^\prime$  and $B\leq N^\prime$.  By  Theorem~\ref{May_26_2024_4}, $f^\prime$ belongs to a back-and-forth system between $M^\prime$ and $N^\prime$ and so it is a partial elementary map from $M^\prime$ to $N^\prime$ and also from $M$ to $N$.
		
		Consider now the general case.  Fix an HF-ordering of $M$ having $A$ as an initial segment.  Consider a finite set $A_0\subseteq A$ and let $A_0^\prime =\cl_<(A_0)$.  Since $A$ is an initial segment, $A_0^\prime\subseteq A$. By Lemma~\ref{Feb_23_2024_7},  we may assume that $A_0^\prime$ is an initial segment of  $A$. Then $A_0^\prime\leq A$  and  $A_0^\prime\leq M$. Now order $B$ in such a way that $f$ is an order isomorphism. This gives an HF-ordering of $B$.  Moreover, if $B_0=f(A_0)$ and $B_0^\prime= f(A_0^\prime)$, then  $B_0\subseteq B_0^\prime$ and $B_0^\prime$ is an initial segment of $B$. Hence  $B_0^\prime\leq B$  and  $B_0^\prime\leq N$. Since $A_0^\prime$ and $B_0^\prime$ are countable, by the initial case we know that $f\restriction A_0^\prime$ is elementary. Hence, $f$ is elementary.
	\end{proof}

\end{document}
